\documentclass[conference]{IEEEtran}
\IEEEoverridecommandlockouts
\usepackage[utf8]{inputenc}
\usepackage[T1]{fontenc}
\usepackage{cite}
\usepackage{amsmath,amssymb}
\usepackage{graphicx}
\graphicspath{{figures/}}
\usepackage{textcomp}
\usepackage{booktabs}
\usepackage[hidelinks]{hyperref}

% IEEEtran yazdirir bunlari Ingilizce; Turkce bildiride Turkce olmali.
\renewcommand{\abstractname}{Öz}
\renewcommand{\IEEEkeywordsname}{Anahtar Kelimeler}
\renewcommand{\figurename}{Şekil}
\renewcommand{\tablename}{Tablo}
\renewcommand{\refname}{Kaynaklar}

\begin{document}

\title{Ne Zaman Bilmediğini Bilmek:\\
GNSS'siz İHA Uçuşu için Sıralı Harita Çapalı\\
Görsel Konumlandırma}

% CIFT KOR DEGERLENDIRME: yazar bilgisi kabul sonrasi eklenecek.
\author{\IEEEauthorblockN{}\IEEEauthorblockA{}}

\maketitle

\begin{abstract}
GNSS reddedildiğinde mutlak konuma ulaşmanın standart yolu, aşağı bakan
kamerayı önceden yüklenmiş bir uydu haritasıyla eşleştirmektir. Güncel
literatür bunu tek görüntülü bir getirme problemi olarak ele alıyor ve yaygın
kıyas kümesi doygunluğa yaklaşmış durumda. Oysa uçan bir hava aracı bir
\emph{dizi} üretir ve bu bağlamda belirleyici kısıt getirme doğruluğu değil,
karelerin önemli bir bölümünün hiç eşleştirilememesidir. Bu çalışmada çapraz
görüş eşleştirmesini görsel odometriyle bir parçacık süzgecinde birleştiriyor
ve 406--2572\,m irtifa, uçuş başına 9--103\,km ve 2016--2023 arası çekim
tarihleri kapsayan \emph{on} gerçek İHA tarama uçuşunda değerlendiriyoruz.
Görüntüsü harita ile kullanılabilir düzeyde örtüşen yedi uçuşta sistem
8,35--24,79\,m medyan hatayı karelerin $\geq$\%99,7'sinde koruyor; buna karşılık
tek başına görsel odometri 74\,km'de 2803\,m sürükleniyor. Aynı düzenekten iki
yan ürün çıkıyor. Birincisi, harita eşleşmelerinin artık dönmesi platformun
yönelim hatasını açığa vuruyor ve sistem manyetometre sapmasını GNSS olmadan
kalibre ederek sürüklenmeyi kat edilen yolun \%3,795'inden \%0,865'ine
indiriyor. İkincisi, planlanan bir rota üzerinde uçmadan ölçülebilen
\emph{eşleşme oranı}, sistemin orada çalışıp çalışmayacağını önceden
söylüyor (logaritmik hatayla $\rho=-0{,}764$, eşik \%50 civarı). Son olarak
genellenebileceğine inandığımız bir olumsuz sonuç bildiriyoruz: LoFTR yerine
RoMa kullanmak çöken uçuşlarda eşleşme oranını \%12--42'den \%96--100'e
çıkarıyor ama uçtan uca doğruluğu \emph{çarpıcı biçimde kötüleştiriyor},
çünkü RoMa ilgisiz görüntüler için de kendinden emin karşılıklar üretiyor.
Konumlandırmada bir eşleyicinin değeri ne kadar eşleştirdiğiyle değil,
susabilmesiyle ölçülür.
\end{abstract}

\begin{IEEEkeywords}
GNSS'siz seyrüsefer, çapraz görüş coğrafi konumlandırma, parçacık süzgeci,
görsel odometri, insansız hava aracı, görüntü eşleştirme
\end{IEEEkeywords}

\section{Giriş}

GNSS alıcısı karıştırılan ya da aldatılan bir insansız hava aracı mutlak
konumunu yitirir. Ataletsel ve görsel odometri bağıl hareketi vermeye devam
eder, ama hataları sınırsız birikir: burada incelenen 74\,km'lik uçuşta kareler
arası görsel odometri gerçek konumdan 2803\,m uzakta bitiyor. Yerleşik çözüm,
coğrafi referanslı bir uydu haritası taşımak ve kameranın gördüğünü bu haritayla
eşleştirerek mutlak konumu geri kazanmaktır.

Güncel çalışmalar problemi çapraz görüş görüntü getirme olarak kuruyor ve
yaygın kıyas kümesini doygunluğa yaklaştırdı; University-1652 bugün \%97,5
Recall@1 seviyesinde~\cite{mgs2net}. Bu kuruluş, işin zor kısmının aday
konumlar arasından doğru olanı seçmek olduğunu sessizce varsayar. Gerçek arazi
üzerindeki uçuşta ise bambaşka bir kısıt baskın çıkıyor: karelerin önemli bir
bölümünde seçilecek \emph{doğru bir eşleşme zaten yok}. Su üstünde, tekdüze
tarlada, tekrar eden yapı desenlerinde ve haritayla kamera mevsim ya da yıl
olarak ayrıştığında eşleştirme kullanılabilir hiçbir şey döndürmüyor. Tek kare
taban çizgimizde bu oran, kümemizdeki en kolay uçuşta \%23, en zorunda \%96.

Uçan bir hava aracı bağımsız fotoğrafların toplamı değildir. Bir yörünge izler
ve eşleştirilemeyen kareler, eşleştirilebilenlerin arasında kalır. Bu yüzden
problemi sıralı durum kestirimi olarak ele alıyoruz: görsel odometri hareket
modelini, harita eşleştirmesi mutlak düzeltmeleri sağlıyor ve başarısız olmasına
izin veriliyor, parçacık süzgeci ikisini birlikte taşıyor. Bu, yer robotiğinde
yerleşik bir uygulama olmakla birlikte güncel İHA çapraz görüş literatüründe
neredeyse hiç yok.

Katkılarımız şunlardır: (i) bir değil \emph{on} gerçek uçuşta değerlendirilen,
her uçuşun kendi ilk \%20'sinden kalibre edilip kalanında ölçülen tam otomatik
bir sıralı harita çapalı konumlandırma sistemi; (ii) harita eşleşmelerinin
geometrisinden, GNSS gerektirmeden yapılan çevrimiçi yönelim ve ölçek
öz-kalibrasyonu; (iii) uçmadan ölçülebilen bir konuşlandırılabilirlik
yordayıcısı; (iv) mekanizması tespit edilmiş bir olumsuz eşleyici sonucu.

\section{İlgili Çalışmalar}

\textbf{Çapraz görüş coğrafi konumlandırma.} Baskın kuruluş, tek bir İHA
görüntüsünü uydu karolarından oluşan bir veri tabanıyla eşleştirir. Güncel
çalışmalar dönüştürücü ve durum uzayı mimarilerini~\cite{mgs2net}, dik olmayan
bakışa dayanıklılığı~\cite{offnadir} ve bozulmuş görüntüyü~\cite{relate}
kapsıyor. 2024--2026 arasındaki yirmi böyle çalışmayı tarayıp kare başına
sonucu zaman içindeki bir hareket modeliyle birleştiren hiçbirine rastlamadık.

\textbf{Hava görüntüsüne karşı sıralı konumlandırma.} Hava haritalarına karşı
parçacık süzgeci \emph{yer} platformlarında yerleşiktir: BEV-Patch-PF öğrenilmiş
kuşbakışı eşleştirmeyi arazi araçları için bir parçacık süzgeciyle
birleştiriyor~\cite{bevpatchpf}. Hava araçları için en yakın çalışma modern
öğrenilmiş eşleştirmeden öncesine dayanıyor: Shan ve ark.~\cite{shan} optik
akışı harita korelasyonuyla süzüyor. Katkımız tam bu kesişimde duruyor: bir
uçakta, sıralı bir süzgecin içinde modern dedektörsüz bir eşleyici ve
kilometrelerce yörünge üzerinde metrik hata.

\textbf{Görüntü eşleştirme.} LoFTR~\cite{loftr} dönüştürücülerle dedektörsüz
kabadan inceye eşleştirme yapar. RoMa~\cite{roma} temel model kodlayıcısıyla
yoğun bir çarpıtma kestirir ve güncel kıyas kümelerinde önde gider. İkisini de
kullanıyor, kıyas sıralamasının bu göreve neden aktarılmadığını
Bölüm~\ref{sec:esleyici}'de bildiriyoruz.

\section{Yöntem}

Her İHA karesi platform yöneliminden kuzey-yukarı döndürülüyor ve haritanın yer
ölçeğine örnekleniyor; sonuç, $300 \times 300$\,m kapsayan metrik ve kuzey
hizalı bir kırpma. Yönelim ataletsel birimden, irtifa barometrik
yükseklikölçerden geliyor; hiçbiri GNSS'e bağlı değil. Harita EPSG:4326'da
saklanıyor; pikselleri derecede kare, metrede değil --- enlemimizde doğu-batı
0,2526\,m/px'e karşı kuzey-güney 0,2974\,m/px, yani \%17,8 anizotropi. Bu
düzeltilmediğinde eşleştirmeye sistematik bir ölçek hatası olarak giriyor.

Mutlak düzeltme için, kestirilen konumun etrafındaki karo DINOv2 tanımlayıcıları
ile getiriliyor ve LoFTR karşılıkları üzerinden RANSAC ile bir benzerlik
dönüşümü kestiriliyor. Yeterli iç nokta bulunamazsa kare \emph{reddediliyor};
sistem yanlış bir konum bildirmek yerine hiçbir şey bildirmiyor. Görüntü
merkezinden hava aracı konumuna geçerken duruş kaynaklı kayma modelleniyor:
$d = h\tan\theta$ ile eğim ve yalpa, gövde çerçevesinde iz-boyu ve iz-dışı
bileşenlere ayrılıyor.

Sıralı füzyonda parçacık süzgeci, hareket modelini görsel odometriden alıyor;
harita düzeltmeleri olabilirlik olarak giriyor ve gelmedikleri karelerde süzgeç
yalnızca hareket modeliyle ilerliyor. Ölçüm güdümlü enjeksiyon, her adımda
parçacıkların bir bölümünü hareket modeli yerine ölçüm dağılımından çekiyor; bu
olmadan süzgeç erken bir yanlış eşleşmeye kilitlenip doğru ölçümü görse bile
kıpırdayamıyor.

\textbf{Haritadan öz-kalibrasyon.} Kuzey-yukarı bir kareyi kuzey-yukarı bir
karoya oturtan dönüşümün artık dönmesi sıfır değilse, yönelim açısının kendisi
tam o kadar yanlıştır. Uçuş boyunca ölçüldüğünde bu hatanın uçuş yönüne göre
değiştiği görülüyor: tarama deseninin iki kolunda $-1{,}93^\circ$ ve
$-7{,}08^\circ$. Yöne bağlı bu imza manyetometrenin sert-demir hatasının klasik
parmak izidir ve GNSS'ten hesaplanan değerlerle ($+1{,}65^\circ$,
$+7{,}22^\circ$) büyüklük ve işaret olarak uyuşuyor. Düzeltmeyi odometriye geri
beslemek sürüklenmeyi kat edilen yolun \%3,795'inden \%0,865'ine indiriyor.

\section{Deneysel Kurulum}

UAV-VisLoc~\cite{uavvisloc} veri kümesini kullanıyoruz; gerçek İHA uçuşlarını
coğrafi referanslı haritalar ve işlenmiş GNSS ile eşleştiriyor. Referans
kalitesini bağımsız olarak doğruladık: düz tarama kollarında iz-dışı saçılma
1,50\,m, yani referansın kendisi bildirdiğimiz hatalara göre temiz.

On bir uçuşun onunu değerlendiriyoruz; dışlanan uçuşun üstverisinde yöntemin
gerektirdiği duruş ve yönelim sütunları yok ve yükleyicimiz bunu desteklenmeyen
sayılar üretmek yerine açıkça reddediyor. Bir uçuşun haritası tek dosya yerine
toplam 4,5\,GB'lık $2\times2$ GeoTIFF ızgarası hâlinde geliyor; bunları hiçbir
kopya çıkarılmadan, GDAL sanal rasterı üzerinden sunuyoruz --- konuşlandırmada
karolu haritalar da zaten böyle gelir. Her uçuşun yer örnekleme aralığı ve
kamera montaj açısı o uçuşun ilk \%20'sinden kestiriliyor, bildirilen bütün
değerler kalan \%80'den geliyor. Bütün deneyler 4\,GB bellekli bir dizüstü
NVIDIA RTX 3050 Ti üzerinde koşuyor; en yüksek bellek kullanımı 1,4\,GB.

\section{Sonuçlar}
\label{sec:sonuclar}

\begin{figure*}[t]
\centering
\includegraphics[width=\textwidth]{fig1_trajectory.pdf}
\caption{Uçuş 03 (768 kare, 74\,km, 466\,m irtifa), 8,8$\times$7,3\,km'lik
ortofotosu üzerinde. İki panel aynı çerçeveyi paylaşıyor. (a) Tek başına görsel
odometri yönelim ve ölçek hatasını biriktirerek haritanın tamamen dışına
çıkıyor ve gerçek konumdan 2803\,m uzakta bitiyor ($\times$). (b) Aynı odometri
harita eşleşmeleriyle çapalandığında her karede gerçek yörüngenin üzerinde
kalıyor; işaretli kareler eşleştirmenin tamamen çöktüğü ve süzgecin yeniden
çapa atana dek yalnızca hareket modeliyle ilerlediği dört kesim.}
\label{fig:yorunge}
\end{figure*}

\begin{table}[t]
\caption{Uçuş başına sonuçlar, tam otomatik. Çizginin üstündeki yedi uçuşun
eşleşme oranı en az \%50; altındaki üçünün değil.}
\label{tab:ucuslar}
\centering
\small
\begin{tabular}{lrrrrr}
\toprule
Uçuş & İrtifa (m) & Yol (km) & Eşleşme & Kapsama & Medyan (m) \\
\midrule
03 & 466  & 74  & \%93 & \%100,0 & \textbf{8,35}  \\
09 & 546  & 77  & \%70 & \%100,0 & \textbf{14,94} \\
06 & 834  & 24  & \%76 & \%99,7  & \textbf{15,06} \\
04 & 544  & 83  & \%90 & \%100,0 & \textbf{15,44} \\
05 & 2313 & 30  & \%50 & \%99,8  & \textbf{16,69} \\
01 & 406  & 66  & \%77 & \%100,0 & \textbf{22,51} \\
11 & 2572 & 84  & \%90 & \%99,8  & \textbf{24,79} \\
\midrule
02 & 406  & 86  & \%37 & \%99,7  & 53,45  \\
10 & 773  & 9   & \%13 & \%84,7  & 126,88 \\
08 & 551  & 103 & \%33 & \%79,7  & 648,49 \\
\bottomrule
\end{tabular}
\end{table}

\begin{figure}[t]
\centering
\includegraphics[width=\columnwidth]{fig3_law.pdf}
\caption{On uçuşta eşleşme oranına karşı nihai doğruluk. Eşleşme oranı gerçek
konum verilerek ölçülüyor, yani algoritmayı değil veriyi betimliyor ve
planlanan bir rota üzerinde uçmadan hesaplanabiliyor. Renk irtifayı gösteriyor
ve grupları ayırmıyor: 2572\,m'deki uçuş (11) çalışıyor, 551\,m'deki (08)
çöküyor.}
\label{fig:yasa}
\end{figure}

Tablo~\ref{tab:ucuslar} uçuş başına sonuçları veriyor. Yedi uçuş 8,35 ile
24,79\,m arasında medyan hatayla ve \%99,7 veya üstü kapsamayla konumlanıyor;
üçü çöküyor.

\textbf{Bileşen karşılaştırması.} Uçuş 03'te üç yaklaşımı doğrudan
karşılaştırdığımızda (Şekil~\ref{fig:yorunge}): tek başına görsel odometri her
kareyi kapsıyor ama 2803\,m sürükleniyor; tek kare harita eşleştirmesi
çözebildiği karelerin \%76,6'sında 5,46\,m medyan tutturuyor ve kare başına
3,80 eşleştirme çağrısı harcıyor; füzyon karelerin \%100'ünü 6,20\,m medyanla
ve kare başına 1,52 çağrıyla kapsıyor. Tek kare taban çizgisi 100\,m'nin
ötesinde kaba aykırı değer üretmiyor: getirme yanlış bölgeyi önerdiğinde
doğrulama onu reddediyor ve sistem yanlış bir konum bildirmek yerine hiçbir şey
bildirmiyor.

\textbf{Çalışan uçuşları çökenlerden ne ayırıyor?} Başarım uçuşlar arasında
78 kat değişiyor; bu, sistemin bazı uçuşlarda mı başarısız olduğu yoksa bazı
uçuşların mı eşleştirilemez olduğu sorusunu doğuruyor. İkincisini doğrudan
ölçüyoruz: gerçek konum verildiğinde karelerin ne kadarı haritayla eşleşiyor?
Arama içermediği için bu, verinin bir özelliği.

Cevap iki grubu ayırıyor (Şekil~\ref{fig:yasa}). Eşleşme oranı $\geq$\%50 olan
uçuşlar $\geq$\%99,7 kapsamada 8,35--24,79\,m medyan veriyor; altındakiler
\%80--85 kapsamada 53--648\,m. Eşleşme oranı ile logaritmik hata arasındaki
korelasyon $-0{,}764$. Ayırt edici olan irtifa değil: 2572\,m'deki uçuş
çalışıyor, 551\,m'deki çöküyor. Bunun pratik değeri, eşleşme oranının uçuş
gerektirmemesinde: planlanan bir rota ve bir harita verildiğinde önceden
ölçülebiliyor ve konuşlandırılabilirlik bir sorudan bir ölçüme dönüşüyor.

\section{Bir Eşleyiciyi Burada Değerli Kılan Nedir}
\label{sec:esleyici}

Yukarıdaki çözümleme darboğazın eşleştirme kalitesi olduğunu gösterdiği için
mantıklı adım LoFTR'ı daha güçlü bir eşleyiciyle değiştirmekti. RoMa güncel
eşleştirme kıyaslarında önde gidiyor. Bilinen doğru konumlarda ölçüldüğünde
sorunu çözmüş görünüyor: çöken üç uçuşta eşleşme oranı \%12--42'den
\%96--100'e çıkıyor, hem de 2,7\,GB bellekle.

Uçtan uca sonuç ise çok daha kötü: uçuş 01'de LoFTR'ın 22\,m verdiği yerde
1669\,m.

Kıyaslama yanlış soruyu soruyordu. Her iki eşleyiciye de yalnızca \emph{doğru}
karoyu gösteriyordu; oysa bir konumlandırma sistemi mesaisinin çoğunu adayın
doğru yer olup olmadığına karar vermeye harcar. Her kareyi haritanın rastgele
ve ilgisiz bir parçasıyla eşleştirdiğimizde:

\begin{center}
\small
\begin{tabular}{lrrr}
\toprule
& Doğru yerde & Yanlış yerde & Oran \\
\midrule
LoFTR & 709 iç nokta & \textbf{0} & 709$\times$ \\
RoMa  & 4598 iç nokta & 341 & 13,5$\times$ \\
\bottomrule
\end{tabular}
\end{center}

LoFTR ilgisiz görüntüde hiçbir şey döndürmüyor; RoMa yüzlerce karşılık
döndürüyor. Görünürdeki üstünlüğü, yalnızca kolay yönü ölçen bir ölçütün yan
ürünüymüş. Kabul eşiğini yeniden ayarlamak da kurtarmıyor: ham iç nokta, iç
nokta oranı ve eşleyici güveni ölçütlerinde LoFTR için sıfır hatalı bir eşik
\emph{var}, RoMa için hiçbir ölçütte yok --- bazı yanlış yerler bazı doğru
yerlerden yüksek puan alıyor.

Uçtan uca, eşikler RoMa lehine ayarlıyken uçuş 08 maliyeti somutlaştırıyor:
LoFTR karelerin \%4,5'inde konum üretiyor, medyan hata 582\,m --- yani
bilmediğini söylüyor. RoMa \%95,5'inde üretiyor ve ortalama 6250\,m yanılıyor.
Uçuşta birinci davranış ataletsel seyre devri tetikler; ikincisi hava aracını
altı kilometre yanlış yere, hiçbir uyarı vermeden uçurur.

Bu yüzden LoFTR'ı koruyoruz; daha güçlü eşleyici olduğu için değil --- değil
--- bu görev eşleştirme kıyaslarının ölçmediği bir özellik istediği için. Bir
eşleyici konumlandırmada, hiçbir şey döndürmeye istekli olduğu ölçüde
değerlidir.

\section{Bileşen Katkıları ve Dayanıklılık}

Bileşenleri teker teker çıkardığımızda (300 kare, sabit tohum) tek en büyük
katkı montaj açısı düzeltmesinden geliyor: onsuz medyan 17,13\,m, onunla
6,62\,m. Odometriyi çıkarmak medyanı 9,98\,m'ye taşıyıp bırakıyor ama kuyruğu
yok ediyor --- \%90 dilim 16,17\,m'den 367,66\,m'ye çıkıyor, ki bir hareket
modeli tam olarak bunun içindir. Ölçüm güdümlü enjeksiyonu çıkarmak 1,0\,m
medyan ve 3,3\,m \%90 dilim maliyetli. İki bileşen ise yerini hak etmiyor ve
bunu aksini iddia etmek yerine bildiriyoruz: parçacık sayısı neredeyse etkisiz
(100 parçacıkla 7,08\,m, 2000 ile 6,80\,m) ve olabilirlik aykırı değer tabanını
kaldırmak ölçülebilir bir şey değiştirmiyor, çünkü kapı ve enjeksiyon kötü
düzeltmeleri zaten daha yukarıda süzüyor.

Yirmi bozulma koşulunda ortaya çıkan örüntü şu: sistem \emph{ışıkölçüme}
kayıtsız, \emph{dokuya} duyarlı. En yüksek şiddette sis 0,6\,m, karenin üçte
birini kapatan bir örtme 0,6\,m, çözünürlük kaybı 0,9\,m maliyetli. Asıl
başarısızlık kipi hareket bulanıklığı: orta düzeyde medyan 6,62'den 58,31\,m'ye,
ağır düzeyde 500\,m'ye çıkıyor. Aynı gözlem bir kısmi çare de veriyor --- harita
karosu zaten keskin ve aynı yeri aynı ölçekte gösterdiği için, karo ile sorgu
arasındaki keskinlik farkı bulanıklığı ölçüyor; karoyu sorguya eşitleyecek kadar
bulandırmak sürdürülmüş bulanıklıkta medyanı 44,47\,m'den 20,00\,m'ye,
\%90 dilimi 214\,m'den 91\,m'ye indiriyor. Temiz başarımı geri getirmiyor;
ağır titreşim gerçek bir sınır olarak duruyor.

Uydu düzeltmelerini rastgele atmak, seyrekleşene kadar az şeye mal oluyor:
kesinti yokken 6,62\,m, düzeltmelerin yarısı atıldığında 7,38\,m, \%75'i
atıldığında 10,76\,m ve \%90'ı atıldığında 41,55\,m.

\section{Sınırlar ve Sonuç}

On uçuşun tamamı tek bir veri kümesinden ve tek bir ülkeden geliyor; irtifa
406--2572\,m ve tarihler 2016--2023 arasında yayılsa da çekim sistemi ve harita
sınıfı hepsinde ortak. Onun üçü çöküyor ve sebebi ölçülmüş olsa da bu gerçek
bir işletme sınırı. İşlem çevrimdışı, dizüstü kartında kare başına 866\,ms; bu
veride kareler 7\,s'de bir geldiği için rahat gerçek zamanlı, ama gömülü
donanımda denenmedi. Düzlemsel homografi varsayımı dike yakın karelerde geçerli,
eğik bakışta bozuluyor; artık hatanın baskın kaynağı bina kabartma kayması ve
ulaşılabilir tabanın 1\,m değil 6\,m civarında olmasının sebebi bu.

GNSS'siz görsel konumlandırmayı kare başına getirme yerine sıralı kestirim
olarak ele almak, hangi problemin zor olduğunu değiştiriyor. Kısıt getirme
doğruluğu değil; kısıt, karelerin çoğunun hiç eşleştirilememesi ve bir hareket
modelinin tam olarak onları kapatması. On gerçek uçuşta füzyon, harita ile
kameranın tanınabilir ölçüde aynı dünyayı gösterdiği her yerde karelerin
\%99,7 veya fazlasını 8--25\,m medyanla konumlandırıyor; bu koşulu sağlayan
kare oranı ise uçuştan önce ölçülebiliyor. Aynı eşleşme geometrisi platformun
kendi yönelim hatasını da veriyor ve manyetometrenin hiçbir konum bilgisi
olmadan haritaya karşı kalibre edilmesini sağlıyor. Eşleyici çalışması ise
eşleştirme literatürünün hâlihazırda bildirmediği bir ölçüt öneriyor:
konumlandırmada seçicilik duyarlılığa baskındır ve her zaman cevap veren bir
eşleyici, bazen susandan kötüdür.

Kod, veri hazırlığı ve bildirilen bütün ölçümler kabul hâlinde herkese açık
olarak yayımlanacaktır.

\begin{thebibliography}{00}
\bibitem{mgs2net} M. Li \textit{et al.}, ``(MGS)$^2$-Net: Macro/micro-geometric
structure and scale modeling for cross-view geo-localization,''
arXiv:2602.10704, 2026.
\bibitem{offnadir} Q. Qiao \textit{et al.}, ``OffNadirLoc: A benchmark for large
off-nadir view geo-localization,'' arXiv:2607.19951, 2026.
\bibitem{relate} H. Jiang \textit{et al.}, ``ReLATE: Reliability-guided evidence
fusion for degraded UAV--satellite matching,'' arXiv:2607.25524, 2026.
\bibitem{bevpatchpf} D. Lee \textit{et al.}, ``BEV-Patch-PF: Particle filtering
with learned aerial patch matching for GPS-free localization,''
arXiv:2512.15111, 2025.
\bibitem{shan} M. Shan \textit{et al.}, ``Google map aided visual navigation for
UAVs in GPS-denied environment,'' arXiv:1703.10125, 2017.
\bibitem{loftr} J. Sun \textit{et al.}, ``LoFTR: Detector-free local feature
matching with transformers,'' in \textit{CVPR}, 2021.
\bibitem{roma} J. Edstedt \textit{et al.}, ``RoMa: Robust dense feature
matching,'' in \textit{CVPR}, 2024.
\bibitem{uavvisloc} W. Xu \textit{et al.}, ``UAV-VisLoc: A large-scale dataset
for UAV visual localization,'' arXiv:2405.11936, 2024.
\end{thebibliography}

\end{document}
